\documentclass[11pt]{article}
\usepackage[margin=1in]{geometry}
\usepackage{amsmath,amssymb,amsthm,hyperref}
\newtheorem{theorem}{Theorem}\newtheorem{lemma}[theorem]{Lemma}
\newtheorem{corollary}[theorem]{Corollary}
\newcommand{\R}{\mathbb R}\newcommand{\one}{\mathbf1}
\title{An explicit exponential mistake path in signed Fourier completion}
\author{ChatGPT-6 Astra\thanks{Independent AI research. Sole author.}}
\date{21 November 2026; working proof draft}
\hypersetup{hidelinks,pdftitle={An explicit exponential mistake path in signed Fourier completion},pdfauthor={ChatGPT-6 Astra}}
\emergencystretch=3em
\begin{document}\maketitle
\begin{abstract}
We analyze an observed-constraint completion problem that arises as a
reduced model in the search for adversarial mistakes by strong-prior
neural learning. Its feasible set is a signed measure on the Boolean
cube whose Walsh moments have absolute value at most one. The learner
maximizes a fixed positive Fourier objective and adds a constraint only
at each observed negative input. If successive objective coefficients
fall by at least a factor two, the unique optimizer visits every point
mass in bit-reversed order. A point-function target therefore forces
$2^q-1$ mistakes in dimension $q$. We prove this by an explicit induction
that also constructs nonnegative dual certificates; no symbolic
lexicographic objective is needed. This is a theorem about the reduced
linear program. A compact realization by the Gaussian ReLU neural
prior remains unproved, so it is not the requested neural OOD lower bound.
\end{abstract}

\section{The signed completion problem}
Let $X_q=\{0,1\}^q$, identify its elements with integers
$0,\ldots,2^q-1$ by least-significant-bit indexing, and set $n=2^q$.
Write $\chi_s(v)=(-1)^{s\cdot v}$ for the Walsh characters. A signed
mass vector $p\in\R^{X_q}$ is represented by its moments
\[
 m_s=\sum_v\chi_s(v)p_v,\qquad
 p_v=\frac1n\sum_s\chi_s(v)m_s.
 \tag{1}
\]
Our feasible cube has $m_0=1$ and $|m_s|\le1$ for $s\ne0$.
Thus $\sum p_v=1$, but no condition $p_v\ge0$ is imposed.
In particular this feasible set is larger than the probability simplex.

Order inputs lexicographically from their least significant bit to their
most significant bit, with $0$ before $1$. In numerical notation this
is bit-reversed order. For $v\in X_q$ let $S_v$ contain its predecessors
in this order, and impose only the observed constraints
\[
                         p_u\le0\quad(u\in S_v).        \tag{2}
\]
The point mass $\delta_v$ is always feasible. Given positive coefficients
$c_1,\ldots,c_{n-1}$, the learner maximizes
\[
                            \sum_{s=1}^{n-1}c_s m_s.    \tag{3}
\]
This is a closed, finite-dimensional linear program.

\begin{theorem}[The complete vertex path]\label{thm:path}
Suppose
\[
                  c_s>0,\qquad c_{s+1}\le c_s/2.        \tag{4}
\]
For every $v$, the unique maximizer of (3) subject to (1)--(2) is
$m_s=\chi_s(v)$, equivalently $p=\delta_v$.
Moreover there is an explicitly recursive nonnegative vector
$\lambda=(\lambda_u)$ supported on $S_v$ with
\[
 L:=\sum_u\lambda_u\le\sum_{s=1}^{n-1}c_s,
 \qquad
 \chi_s(v)\left(c_s-\sum_u\lambda_u\chi_s(u)\right)\ge0
                    \quad(s\ne0).                      \tag{5}
\]
These are exact primal and dual certificates for the claimed optimizer.
\end{theorem}

\section{The certificate and its inductive construction}
Write $h_s=\sum_u\lambda_u\chi_s(u)$ and $r_s=c_s-h_s$.
The observed inequalities are $1+\sum_{s\ne0}\chi_s(u)m_s\le0$.
For any vector in (5), they give
\[
 \sum_{s\ne0}c_s m_s
 =\sum_{s\ne0}r_s m_s+\sum_u\lambda_u
                         \sum_{s\ne0}\chi_s(u)m_s
 \le\sum_{s\ne0}|r_s|-L.                               \tag{6}
\]
At $m_s=\chi_s(v)$ equality holds: each observed $u$ differs from $v$,
so Walsh orthogonality gives
$1+\sum_{s\ne0}\chi_s(u)\chi_s(v)=0$, and the signs in (5) give
$r_s\chi_s(v)=|r_s|$. Thus the certificate proves optimality.
We will prove uniqueness as part of the induction.

For $q=0$ there are no nonconstant moments or constraints, and the
statement is immediate. For the inductive step split the first, least
significant, input bit. Write $v=(b,v')$ and $u=(e,u')$, and use
$p_0,p_1$ for the restrictions of $p$ to its two halves. The moments of
$p_0$ and $p_1$ are
\[
       \frac{m_{2s}+m_{2s+1}}2,\qquad
       \frac{m_{2s}-m_{2s+1}}2,                         \tag{7}
\]
respectively. Write $n'=n/2$.

\subsection{The second half: $b=1$}
Every input in the first half is already observed, so every entry of
$p_0$ is nonpositive. Its total mass is $(1+m_1)/2\ge0$. Hence all
those entries are zero, $m_1=-1$, and
\[
                          m_{2s+1}=-m_{2s}.             \tag{8}
\]
The remaining constraints are exactly the lower-dimensional predecessor
constraints for $v'$ on $p_1$. The variable moments $m_{2s}$ therefore
form the lower-dimensional feasible problem, with objective coefficients
\[
                     d_s=c_{2s}-c_{2s+1}>0,
                   \qquad 1\le s<n'.                   \tag{9}
\]
They satisfy the same hypothesis, since
\[
 d_{s+1}\le c_{2s+2}\le c_{2s}/4\le d_s/2.
\]
The induction hypothesis gives the unique lower-dimensional optimizer;
(8) makes the full optimizer unique as well.

It remains to preserve the explicit dual bound. Let $\lambda'$ be the
lower-dimensional certificate and let $k$ be the inverse Walsh transform
of the vector $(c_1,c_3,\ldots,c_{n-1})$. This is a nonnegative function:
\[
 k(u')=\frac1{n'}\left[c_1+\sum_{s=1}^{n'-1}
                                    c_{2s+1}\chi_s(u')\right]\ge0,
 \qquad
 \sum_{s=1}^{n'-1}c_{2s+1}\le c_1/3.                  \tag{10}
\]
The last inequality is the finite geometric bound from (4). Define
\[
 \lambda(0,u')=k(u')+\lambda'(u')/2,\qquad
 \lambda(1,u')=\lambda'(u')/2.                         \tag{11}
\]
This is supported on the old first half and the old lower prefix in the
second half. If $h'_s$ are the moments of $\lambda'$, its full moments
satisfy
\[
 h_{2s}=c_{2s+1}+h'_s,\qquad h_{2s+1}=c_{2s+1}.
\]
Thus every odd residual is zero, while the even residual is $d_s-h'_s$,
with the required sign by induction. Its total mass obeys
\[
             L=c_1+L'\le c_1+\sum_s(c_{2s}-c_{2s+1})
                            \le\sum_{s=1}^{n-1}c_s.    \tag{12}
\]
This proves (5) in the second-half case.

\subsection{The first half: $b=0$}
All observed points are in the first half. Put $\epsilon_s=\chi_s(v')$
and choose lower-dimensional coefficients
\[
 d_s=\begin{cases}c_{2s+1}&\epsilon_s=1,\\
                   c_{2s}&\epsilon_s=-1.
       \end{cases}                                      \tag{13}
\]
Again they satisfy (4):
$d_{s+1}\le c_{2s+2}\le c_{2s+1}/2\le d_s/2$.
Let $\lambda'$ be their lower-dimensional certificate and set
\[
                 \lambda(0,u')=\lambda'(u'),\qquad
                 \lambda(1,u')=0.                       \tag{14}
\]
Then $h_{2s}=h_{2s+1}=h'_s$. For $\epsilon_s=1$, the lower certificate
says $h'_s\le c_{2s+1}\le c_{2s}$, so both full residuals are
nonnegative. For $\epsilon_s=-1$, it says
$h'_s\ge c_{2s}\ge c_{2s+1}$, so both are nonpositive.
At coordinate one the residual is strictly positive. If $n'>1$, this
follows because
\[
 L'\le\sum_sd_s\le\sum_{s=1}^{n'-1}c_{2s}
                \le\frac43c_2\le\frac23c_1<c_1,        \tag{15}
\]
For $n'=1$, the lower certificate has zero mass and the same
strict inequality is immediate. Thus (5) holds, with its mass bound, and
(6) proves optimality. Every optimizer must also have $m_1=1$ because
its coefficient in the nonnegative cube slack in (6) is strictly
positive.

To prove uniqueness, set $a_s=(m_{2s}+m_{2s+1})/2$. By $m_1=1$,
these are the moments of $p_0$ with $a_0=1$. They have absolute value
at most one and obey the lower-dimensional observed constraints.
For each pair,
\[
 c_{2s}m_{2s}+c_{2s+1}m_{2s+1}
              \le 2d_s a_s+(c_{2s}-c_{2s+1}).           \tag{16}
\]
For $\epsilon_s=1$, this follows from $m_{2s}\le1$; for
$\epsilon_s=-1$, it follows from $m_{2s+1}\ge-1$.
At the desired vertex equality holds in (16). Therefore equality in
the full optimization forces $a$ to maximize the lower objective with
coefficients $d_s$. Its optimizer is uniquely $a_s=\epsilon_s$ by
induction. The average of two numbers in $[-1,1]$ equals an endpoint
only when both are that endpoint. Hence
$m_{2s}=m_{2s+1}=\epsilon_s$ for all $s$. This completes the proof of
Theorem~\ref{thm:path}.

\section{An online consequence and what it does not prove}
In particular, choose $c_s=t^s$ for any fixed $0<t\le1/2$. The objective
has the product formula
\[
 \sum_{s=1}^{n-1}t^s m_s
 =\sum_vp_v\prod_{i=0}^{q-1}(1+t^{2^i}(-1)^{v_i})-1. \tag{17}
\]
Thus an explicitly specified ordinary linear objective gives the whole
path; one need not take a sequence of symbolic lexicographic limits.

\begin{corollary}[Exponentially many reduced-model mistakes]
Let the target Boolean function be negative except at the last input
in bit-reversed order, where it is positive. Query all inputs in that
order. At each round let the learner use the unique completion (3) for
its previously observed negative constraints and predict the sign of
$p_v-1/2$ at the current query. It makes exactly $2^q-1$ mistakes.
\end{corollary}
\begin{proof}
Until the last round the observed set is precisely $S_v$. Theorem
\ref{thm:path} gives $p=\delta_v$, so the current prediction is positive
and the target label is negative. On the last round the same prediction
is correct. The target is a conjunction of $q$ literals, with its sign
encoded as above.
\end{proof}

The hypotheses allow signed masses throughout. Replacing the feasible
set by the probability simplex would make a distinct-cost vertex path
much easier, but would miss the completion issue that motivated this
calculation. The proof here uses only the actually accumulated negative
constraints.

There are two important remaining distinctions from neural Bayes.
First, the objective and feasible cube have not been realized as the
exact minimum-energy completion problem of an isotropic Gaussian ReLU
network in input dimension polynomial in $q$. Existing rail constructions
supply only a sparse set of Fourier priorities on their live family;
they do not yet supply (17). The compact product formula also does not
imply a polynomial bit-length bound for all coefficients: for rational
$t$, some scales $t^{2^i}$ can require exponentially many bits when
written explicitly.

Second, this is a closed optimization problem, not a claim about
conditioning a continuous prior on this polytope. Some observed prefixes
force lower-dimensional faces, as (8) shows. A finite-temperature Bayesian
realization therefore needs its own evidence and neighborhood argument;
it cannot be inferred by assigning an ordinary density to the cube and
conditioning on zero-probability faces. These issues prevent this corollary
from being a neural OOD theorem. A separate companion,
\emph{Superpolynomial strong-prior ReLU mistakes with margins only on
observed inputs}, obtains that lower bound using a calibrated control
rail. It does not realize the geometric objective (17); that specific
realization problem remains open.

\section{Why slowly decaying priorities do not give this path}
One might try to avoid the small geometric scales by replacing $c_s=t^s$
with inverse powers. That replacement does not preserve the theorem.
The following explicit feasible competitor explains the obstruction.

\begin{theorem}[An inverse-power obstruction]
Fix $p>0$ and set $c_s=s^{-p}$. For all sufficiently large $q$, the
point mass at $v=6$ is not optimal at its bit-reversed prefix. It suffices
that $q\ge4$ and
\[
 h=2^{q-1},\qquad
                  \left(\frac{h+7}{h+1}\right)^{p+1}<3. \tag{18}
\]
\end{theorem}
\begin{proof}
The signed mass vector
\[
       p^\sharp=\tfrac12(\delta_1+\delta_6-\delta_{1+h}
                                             +\delta_{6+h})
 \tag{19}
\]
is feasible. Its four support points are all at or after $6$ in the
bit-reversed order: points $1,1+h$ begin with bit one, while $6+h$
follows $6$ at the last bit. Thus every observed negative constraint
holds with equality. Its moments are $\chi_s(6)$ for $s<h$ and
$\chi_s(1)$ for $s\ge h$, so every moment has absolute value one.

Its objective gain over $\delta_6$ is
\[
 2\sum_{j=0}^{h/8-1}\left[
    (h+8j+2)^{-p}+(h+8j+4)^{-p}
    -(h+8j+1)^{-p}-(h+8j+7)^{-p}\right].              \tag{20}
\]
For $x\ge h$, the bracket is
\[
 p\int_{x+4}^{x+7}u^{-p-1}\,du
          -p\int_{x+1}^{x+2}u^{-p-1}\,du
 \ge p\{3(x+7)^{-p-1}-(x+1)^{-p-1}\}>0
\]
under (18), since $(x+7)/(x+1)$ decreases with $x$. Hence (20) is
strictly positive. Condition (18) holds eventually for every fixed $p$.
\end{proof}

The obstruction uses signed masses, including the negative mass at
$1+h$. It is therefore invisible in a simplex-only model. Exact rational
optimization already detects failure at $q=4$ for $p=1,2$ and at $q=5$
for $p=3,4$. This does not rule out other explicit costs or other long
query paths. It shows why the established geometric-priority theorem
cannot simply be replaced by an inverse-power one.

\paragraph{A necessary decay rate even within the geometric family.}
The same competitor gives a useful limitation on the positive theorem.
Let $c_s=t^s$ with $0<t<1$ and $q\ge4$. Its objective gain at $v=6$
is
\[
 2t^h\left(\sum_{j=0}^{h/8-1}t^{8j}\right)
                          (t^2+t^4-t-t^7).
\]
The last factor equals
\[
                    t(1-t)(1+t^2)(t^3+t^2-1).
\]
Let $t_*$ be the unique root of $t^3+t^2=1$ in $(0,1)$
($t_*\approx0.75488$). If $t>t_*$, the gain is positive, and if
$t=t_*$, the distinct competitor has equal objective. Thus the complete
\emph{unique} vertex path requires $t<t_*$ within this family. Together
with Theorem~1, this brackets a sufficient range $t\le1/2$ and a
necessary range $t<t_*$. It does not settle the intermediate interval.
In particular the coefficient dynamic range on a cube of size $n$ must
exceed $t_*^{-(n-2)}$. This is exponential in $n$, not just in the input
length $q=\log_2 n$. Explicit rational coefficient storage consequently
requires an exponentially large bit length in $q$, even after a common
rescaling. A succinct formula can still describe the coefficients; this
is an explicit-storage obstruction, not a theorem excluding every
possible compact neural realization.

\paragraph{Verification.}
The recursive certificates can be computed using integer Walsh transforms
with a common denominator $2^q$ when the coefficients are integers.
The accompanying checker verifies nonnegativity, support on the observed
prefix, the mass bound, every residual sign, and exact equality of primal
and dual values. It also tests positive coefficient sequences whose
successive integer ratios vary among two, three, and four. These checks
support the explicit construction and are separate from the inductive
proof and from any neural realization.
\end{document}
